% dual-mode-timing.tex  --  Dual-mode retrieval over a 6-week timeline
% \input{} file: requires tikz, xcolor[dvipsnames], calc, arrows.meta,
%   patterns, decorations.pathreplacing loaded in the preamble.

\begin{figure}[htbp]
\centering
\begin{tikzpicture}[
    >=Stealth,
    every node/.style={font=\small},
    week label/.style={font=\footnotesize, text=black!70},
    lane label/.style={font=\small\bfseries, anchor=east},
    event label/.style={font=\scriptsize, text=BrickRed, align=center},
    annotation/.style={font=\scriptsize, align=left},
  ]

  %% === Geometry ===
  %% 6 weeks x 2.0 cm = 12.0 cm timeline + labels
  \def\tlength{12.0}        % total timeline length (cm)
  \def\weekw{2.0}           % width per week
  \def\pullY{2.0}           % pull lane centre
  \def\pushY{0.0}           % push lane centre
  \def\laneH{0.7}           % half-height of lane band
  \def\barW{0.06}           % half-width of a pull bar
  \def\barH{0.55}           % pull bar height
  \def\dotR{0.10}           % push marker radius

  %% === Background lane bands ===
  \fill[MidnightBlue!6]
    (0, \pullY-\laneH) rectangle (\tlength, \pullY+\laneH);
  \fill[OliveGreen!6]
    (0, \pushY-\laneH) rectangle (\tlength, \pushY+\laneH);

  %% === Lane labels ===
  \node[lane label, text=MidnightBlue] at (-0.25, \pullY)
    {Pull-режим};
  \node[lane label, text=OliveGreen]   at (-0.25, \pushY)
    {Push-режим};

  %% === Horizontal timeline axis ===
  \draw[->, thick, black!60]
    (0, -1.3) -- (\tlength+0.3, -1.3);
  \node[font=\footnotesize, anchor=west] at (\tlength+0.35, -1.3) {$t$};

  %% Week tick marks and labels
  \foreach \w in {1,...,6} {
    \pgfmathsetmacro{\wx}{(\w - 0.5)*\weekw}
    \draw[black!40] (\wx, -1.20) -- (\wx, -1.40);
    \node[week label, below] at (\wx, -1.42)
      {Тижд.\,\w};
    %% light vertical gridlines across both lanes
    \draw[black!8, very thin]
      (\wx, \pushY-\laneH) -- (\wx, \pullY+\laneH);
  }

  %% Week boundary lines (dashed)
  \foreach \b in {1,...,5} {
    \pgfmathsetmacro{\bx}{\b*\weekw}
    \draw[black!15, densely dashed, thin]
      (\bx, \pushY-\laneH) -- (\bx, \pullY+\laneH);
  }

  %% ---------------------------------------------------------------
  %%  Pull-mode bars  (MidnightBlue)
  %%  Week 1: 0--2, Week 2: 2--4, Week 3: 4--6, Week 4: 6--8,
  %%  Week 5: 8--10, Week 6: 10--12
  %% ---------------------------------------------------------------
  %% Weeks 1--2: dense activity (Operator A)
  \foreach \x in {0.20, 0.50, 0.75, 1.00, 1.25, 1.50, 1.80,
                   2.30, 2.55, 2.85, 3.10, 3.40, 3.65, 3.90} {
    \fill[MidnightBlue, opacity=0.80]
      (\x-\barW, \pullY-\barH/2) rectangle (\x+\barW, \pullY+\barH/2);
  }

  %% Weeks 3--4: sparse/dormant (transition period)
  \foreach \x in {4.70, 5.90, 7.20} {
    \fill[MidnightBlue, opacity=0.35]
      (\x-\barW, \pullY-\barH/2) rectangle (\x+\barW, \pullY+\barH/2);
  }

  %% Weeks 5--6: dense again (Operator C / returned A)
  \foreach \x in {8.50, 8.80, 9.10, 9.35, 9.60, 9.85,
                   10.15, 10.40, 10.65, 10.90, 11.15, 11.45, 11.75} {
    \fill[MidnightBlue, opacity=0.80]
      (\x-\barW, \pullY-\barH/2) rectangle (\x+\barW, \pullY+\barH/2);
  }

  %% ---------------------------------------------------------------
  %%  Push-mode markers  (OliveGreen)
  %% ---------------------------------------------------------------
  %% Weeks 1--2: occasional background refresh
  \foreach \x in {1.00, 3.10} {
    \fill[OliveGreen, opacity=0.30]
      (\x, \pushY) circle (\dotR);
  }
  %% Weeks 3--4: frequent push updates during dormant/transition
  \foreach \x in {4.40, 4.85, 5.30, 5.75, 6.20, 6.65, 7.10, 7.55, 7.95} {
    \fill[OliveGreen, opacity=0.85]
      (\x, \pushY) circle (\dotR);
  }
  %% Weeks 5--6: occasional again
  \foreach \x in {9.60, 11.45} {
    \fill[OliveGreen, opacity=0.30]
      (\x, \pushY) circle (\dotR);
  }

  %% ---------------------------------------------------------------
  %%  Operator rotation markers  (BrickRed dashed verticals)
  %% ---------------------------------------------------------------
  %% Rotation 1: end of week 2 / start of week 3 (x = 4.0)
  \draw[BrickRed, thick, densely dashed]
    (4.15, \pushY-\laneH-0.15) -- (4.15, \pullY+\laneH+0.15);
  \node[event label, above] at (4.15, \pullY+\laneH+0.20)
    {Ротація\\[-1pt]\scriptsize Оператор~A\,$\to$\,B};

  %% Rotation 2: end of week 4 / start of week 5 (x = 8.0)
  \draw[BrickRed, thick, densely dashed]
    (8.25, \pushY-\laneH-0.15) -- (8.25, \pullY+\laneH+0.15);
  \node[event label, above] at (8.25, \pullY+\laneH+0.20)
    {Ротація\\[-1pt]\scriptsize Оператор~B\,$\to$\,C};

  %% ---------------------------------------------------------------
  %%  Phase labels below the timeline
  %% ---------------------------------------------------------------
  %% Active phase (weeks 1--2)
  \draw[decorate, decoration={brace, mirror, amplitude=4pt},
        MidnightBlue!70, thick]
    (0, -2.05) -- (4.10, -2.05);
  \node[font=\scriptsize, text=MidnightBlue!80, below] at (2.05, -2.20)
    {Активна фаза (Оператор~A)};

  %% Dormant / transition phase (weeks 3--4)
  \draw[decorate, decoration={brace, mirror, amplitude=4pt},
        OliveGreen!70, thick]
    (4.20, -2.05) -- (8.20, -2.05);
  \node[font=\scriptsize, text=OliveGreen!80, below] at (6.20, -2.20)
    {Перехідний період --- push-оновлення};

  %% Resumed phase (weeks 5--6)
  \draw[decorate, decoration={brace, mirror, amplitude=4pt},
        MidnightBlue!70, thick]
    (8.30, -2.05) -- (\tlength, -2.05);
  \node[font=\scriptsize, text=MidnightBlue!80, below] at (10.15, -2.20)
    {Відновлення (Оператор~C)};

  %% ---------------------------------------------------------------
  %%  Annotation: pull returns current context at week 3
  %% ---------------------------------------------------------------
  %% Arrow from push lane at ~week 3 up to pull bar at week 3
  \draw[->, OliveGreen!80, thick, shorten >=2pt, shorten <=2pt]
    (4.70, \pushY+\dotR+0.08) -- (4.70, \pullY-\barH/2-0.05);
  \node[annotation, text=OliveGreen!90,
        anchor=south west, inner sep=2pt,
        fill=white, fill opacity=0.85, text opacity=1,
        rounded corners=1.5pt]
    at (4.85, 0.70)
    {Pull-запит Оператора~B\\повертає \emph{актуальний}\\контекст завдяки\\push-оновленням};

  %% ---------------------------------------------------------------
  %%  Legend
  %% ---------------------------------------------------------------
  \node[anchor=north west, inner sep=0pt] at (0, -2.80) {%
    \begin{tikzpicture}[every node/.style={font=\scriptsize, anchor=west}]
      \fill[MidnightBlue, opacity=0.80] (0,0) rectangle (0.3, 0.22);
      \node at (0.42, 0.11) {--- pull-запит (активна сесія)};
      %
      \fill[OliveGreen, opacity=0.85] (4.5, 0.11) circle (0.09);
      \node at (4.7, 0.11) {--- push-оновлення (фон)};
      %
      \draw[BrickRed, thick, densely dashed] (8.5, 0.0) -- (8.5, 0.22);
      \node at (8.65, 0.11) {--- ротація оператора};
    \end{tikzpicture}%
  };

\end{tikzpicture}
\caption{%
  Дворежимне витягування протягом шеститижневого циклу місії.
  У~pull-режимі (верхня доріжка) кожна сесія оператора ініціює
  контекстний запит; у~push-режимі (нижня доріжка) система
  періодично оновлює кешований контекст завдань.
  Під~час ротації оператора (штрихові лінії) push-оновлення
  гарантують, що перший pull-запит нового оператора
  повертає актуальний стан місії без~ручної реініціалізації.%
}
\label{fig:dual-mode}
\end{figure}
